
%=====================================================================
\section{Asymptotics and the Baker obstruction}\label{sec:baker}
%=====================================================================
\subsection{Two-dimensional Euler--Maclaurin}\label{sec:EM}
All radial expansions in this paper rest on the following expansion for cone sums.
\begin{proposition}[Cone Euler--Maclaurin]\label{prop:EM5}
For every rational $x$, the shifted-lattice cone sums of $e^{-tQ}$ that make up $H_s(\e(x)e^{-t})$ and $E_L(x+h)$ have a full asymptotic expansion in powers of
$t^{1/2}$. The coefficients are area terms, edge Bernoulli sums with the half-weight rule $\tilde B_r$ at the edges, and corner terms. The remainders are explicit integrals of periodic Bernoulli functions and are bounded uniformly in every sector $|\arg t|\le\pi/2-\delta$.
\end{proposition}
\catA{Proved, by a one-dimensional lemma (Zagier's asymptotic lemma with the Euler--Maclaurin remainder
$\frac{(-1)^{p+1}}{p!}\int_0^\infty\tilde B_p(y-u)f^{(p)}(y)\,dy$), a quadrant lemma, and a decomposition of the closed cone into translated quadrants.
The sector-uniform bound is order-$3$ two-variable Euler--Maclaurin with every remainder bounded; only $|\tilde B_p|$ enters the bounds. The written
proof is in the companion report (Appendix~\ref{app:repro}); its structure is reproduced here.}
In particular we use the following consequence, which we call \textbf{R1}: for rational $x$, any cone sum $f$ of this type and $|\arg t|\le\pi/2-\delta$,
$|t|\le1$,
\[
f(\e(x)e^{-t})=A/t+h\,t^{-1/2}+\kappa+O(|t|^{1/2}),
\]
with the coefficients of the real-$t$ expansion. Thus radial expansions hold non-tangentially.


\subsection{Radial expansions}
Table~\ref{tab:radial} lists the coefficients of $t^{-1},t^{-1/2},t^0$ at $q=e^{-t}$. The method is two-dimensional
Euler--Maclaurin on shifted quadrants; for $H_s$ the closed cone $\langle(0,1),(21,-1)\rangle$ is decomposed into $21$ residue
classes $(s,0)+\mathbb N v_1+\mathbb N v_2$ minus half-walls (decomposition exact to $q^{2000}$).
\begin{table}[ht]\centering\small
\begin{tabular}{lccc}\toprule
series & $t^{-1}$ & $t^{-1/2}$ & $t^0$\\\midrule
$\Porth_s$ & $I$ & $\sqrt\pi(49-2\sqrt{14})/140=0.5256172043$ & $(11-7I)/20=0.4767030629$\\
$\Cthird$ & $I$ & $-\sqrt\pi(49-2\sqrt{14})/140$ & $(11-7I)/20$\\
$R$ & $I$ & $-0.3606097211$ & $(1-7I)/20$\\
$E$ & $2I$ & $0$ & $(1-7I)/10$\\
$H_s$ & $0$ & $a_H=\sqrt\pi(\sqrt{1190}-17\sqrt{14})/170=-0.3035254127$ & $0$\\
$C$ & $I$ & $a_C=\sqrt\pi(833+204\sqrt{14}-14\sqrt{1190})/2380=0.8291426171$ & $(11-7I)/20$\\\bottomrule
\end{tabular}
\caption{Radial coefficients at $q=e^{-t}$. Status as in Theorem~\ref{thm:C}; the $H_s$ and $C$ entries are the closed forms of Theorem~\ref{thm:Hs}.}\label{tab:radial}
\end{table}
Consistency: $2\Porth_s=E+1+\theta_3-\Thol$ forces the $t^{-1/2}$ coefficient of $\Thol$ to be $\sqrt\pi-2a_P=\sqrt\pi(21+2\sqrt{14})/70$.
For $Z(s)=\sum_{(n,k)\ge0,\ne0}Q_7^{-s}$: simple poles at $s=1$ (residue $I$) and $s=\tfrac12$ (residue $(49-2\sqrt{14})/140$), and
$Z(0)=-(9+7I)/20$; the $\log\varepsilon$ in $Z(0)$ comes entirely from the shift $\tfrac7{20}$.

\begin{lemma}[Zwegers parts carry no $1/t$ term]\label{lem:zw0}
The holomorphic part of any signature-$(1,1)$ indefinite theta $\vartheta^{c_1,c_2}_{a,b}$ with $c_1,c_2$ of negative norm and rational
$a,b$ has $1/t$ coefficient $0$ at $q=e^{-t}$. This covers $H_s$, $\Thol$ and $q^8f_{2,7,7}$.
\end{lemma}
\begin{proof}[Sketch]
The weight is $+1$ on a cone $\mathcal C$ and $-1$ on $-\mathcal C$ (walls contribute $O(t^{-1/2})$). The mean of $\e(B(\nu,b))$ over the coset is the same on
$\mathcal C$ and $-\mathcal C$, and $Q$ is even, so the leading integrals cancel.
\end{proof}
\catA{Proved (symmetry argument; null cones excluded). Numerically $t\,H_s\to0$ (fit $-0.0002$).}

\subsection{The additive obstruction (Theorem~\ref{thm:E}(i))}
Let $\mathcal M$ be the $\Qbar$-span of: (i) $q^rf$, $r\in\Q$, $f$ a holomorphic weight-$1$ form on some $\Gamma(N)$ with $\Qbar$
coefficients; (ii) $q^rf$ with $f$ of half-integral weight, and $f$ of integral weight $k\ge2$ without $q^r$ shift; (iii) $q^r$ times
Zwegers holomorphic parts as in Lemma~\ref{lem:zw0}; (iv) unary, partial and false theta series
$\sum_{n\ge n_0}\psi(n)n^\epsilon q^{an^2+bn}$ with $\psi$ periodic and $\Qbar$-valued, $\epsilon\in\{0,1\}$, $a>0$.
\begin{proof}[Proof of Theorem~\ref{thm:E}(i)]
Multiply by $t$ and let $t\to0^+$. Class (i) contributes $2\pi\beta$ with $\beta\in\Qbar$ (algebraic expansion of $f|S$); (ii) contributes $0$;
(iii) contributes $0$ by Lemma~\ref{lem:zw0}; (iv) contributes $0$ or $\mathrm{mean}(\psi)/(2a)\in\Qbar$. So $I=2\pi\beta+\gamma'$ with
$\beta,\gamma'\in\Qbar$, i.e.\ $\log(6+\sqrt{35})+4\sqrt{35}\,i\beta\log(-1)-2\sqrt{35}\gamma'=0$, a nontrivial $\Qbar$-linear relation
among $1$, $\log(6+\sqrt{35})$, $\log(-1)=i\pi$. Since $\log(6+\sqrt{35})$ (real, nonzero) and $i\pi$ are $\Q$-linearly independent,
Baker's theorem \cite{Baker} says $1,\log\alpha_1,\log\alpha_2$ are $\Qbar$-linearly independent: contradiction.
\end{proof}
\catA{Proved, conditional on Baker's theorem \cite{Baker} and on the algebraicity of cusp expansions (Shimura's $\Q(\zeta_N)$-rationality, i.e.\ the
$q$-expansion principle), both cited. The radial expansion used is R1 (\S\ref{sec:EM}).}
\catB{Reading: the remainder of $C$ is ``weight-1 sized'' at $q\to1$ with a regulator as coefficient; weight-$\tfrac12$ false thetas are
$O(t^{-1/2})$ and cannot carry it. An Appell--Lerch formula for $E$ or $\Porth_s$ with algebraic data would likewise be excluded, but this needs
the $q\to1$ asymptotics of Appell--Lerch sums, which were not verified; that obstruction is \emph{conditional}.}

\subsection{The multiplicative obstruction (Theorem~\ref{thm:E}(ii))}
\begin{proof}
Along $q=e^{-t}$, $F/g=c\,t^{-j}e^{-\alpha/t}(1+O(e^{-\beta/t}))$ times a power series in $t$, so only integer-step powers of $t$ occur;
but $\Cthird$ has the nonzero term $-a_Pt^{-1/2}$ and $C$ the nonzero term $a_Ct^{-1/2}$ (Table~\ref{tab:radial}).
\end{proof}
\catA{Proved. For $C$ the nonvanishing of $a_C$ rests on its closed form (Theorem~\ref{thm:C}). This covers every space with
$k\in\{\tfrac12,1,\tfrac32,2\}$ and $N\mid140$ or $280$; no Sturm-bound search is needed.}

\subsection{Gauss means and the scalar $W_7$ obstruction}\label{sec:gauss}
\begin{proof}[Proof of Theorem~\ref{thm:D} (pattern)]
If $2\mid m$, shift $k\mapsto k+m$: $(Q_7(n,k+m)-Q_7(n,k))/m=7n+7k+\tfrac72(m+1)$ and the phase $\e(7h(m+1)/2)=-1$, so $G=0$. If $m$ is odd,
$Q_7$ is $m$-periodic and $2$ is invertible mod $m$; with $n'=n+7\cdot2^{-1}k$, $hQ_7\equiv h[n'^2-\tfrac{35}4k^2+\tfrac72k]$ and
$G=m^{-2}g_ng_k$ with $|g_n|=\sqrt m$. If $5\mid m$, the shift $(n,k)\mapsto(n-m/5,k+m/5)$ changes the constant part by $j/5$, $j\equiv1\bmod5$,
so $G=0$. If $\gcd(m,35)=1$, completing the square gives $|g_k|=\sqrt m$. If $7^e\Vert m$, $5\nmid m$, by CRT the $7$-part of $g_k$ is $7\sum_{k\bmod 7^{e-1}}\e\bigl(h(-\tfrac54k^2+\tfrac k2)/7^{e-1}\bigr)$, of modulus $7^{(e+1)/2}$,
so $|G|=\sqrt7/m$.
\end{proof}
\catA{Proved; verified on all $3374$ coprime pairs with $m\le105$, zero mismatches. For $x=h/7$, $G=\bigl(\tfrac h7\bigr)i/\sqrt7$.}

For the scalar $W_7$ question: a singularity $\sim G(x)I/(2\pi\varepsilon')$ at $x$ (with $\tau=x+i\varepsilon'$) maps under
$W_7$ to $\sim|7x|^2G(-1/(7x))I/(2\pi\varepsilon')$, so boundedness of $C-\lambda(7\tau)^{-k}C(-1/(7\tau))$ would force
$G(x)=\lambda(7x)^{-k}|7x|^2G(-1/(7x))$ for all $x$. This fails at the $178$ zero-mismatch points; on the $60$ pairs of common support the
moduli obey an exact \emph{weight-1} relation with $|\lambda|=7^{-1/2}$ but the phase varies ($\pi/2$ at $1/7$ and $1$, $-\pi/6$ at $1/3$ and
$3/7$); for $W_{35},W_{70},W_{140}$ there are no pairs with both values nonzero. The map $\tau\mapsto-1/(7\tau)$ does not normalise $\Gamma_0(70)$.
\catA{Zero/nonzero pattern and moduli: exact, by Theorem~\ref{thm:D} (the floating-point exponential sums, nonzero values $\ge0.034$ and zeros $<10^{-9}$, agree).}
\catB{The phase values quoted above are floating-point (with a wide margin) and are not separately proved.}
\catB{The scalar failure is \emph{vector mixing}: the leading profiles obey a vector-valued law (\S\ref{sec:lead}), under which $W_7$-type
maps mix components; this is not a contradiction with modularity of the vector.}
